\documentclass[11pt]{article}
\usepackage[a4paper,margin=25mm]{geometry}
\usepackage{amsmath,amssymb,amsthm,booktabs,hyperref}
\hypersetup{colorlinks=true,urlcolor=blue,linkcolor=blue}
\newcommand{\Q}{\mathbb Q}
\newcommand{\F}{\mathbb F}
\newcommand{\Z}{\mathbb Z}
\newcommand{\divi}{\operatorname{div}}
\newcommand{\Jac}{\operatorname{Jac}}
\newtheorem{theorem}{Theorem}
\title{An infinite family of geometrically simple genus-two Jacobians\\
with rational torsion $\Z/2\Z\times\Z/18\Z$}
\author{}
\date{16 November 2026}
\begin{document}
\maketitle

\noindent
We construct a rational family containing infinitely many geometrically
nonisomorphic genus-two curves whose Jacobians are geometrically simple
and have exact rational torsion $\Z/2\Z\times\Z/18\Z$.
The family contains the known census curve $\texttt{2.18.a}$.

\begin{theorem}\label{thm}
For $k\in\Q$, set
\[
 s=\frac{k^2+3}{4},\qquad
 r=\frac{k^2+3}{(k-1)^2},\qquad
 \lambda=\frac{(s-1)^2}{s^3},
\]
and define
\[
 R(x)=x^3-(x-1)(x-r)(x-s),\qquad
 A(x)=\frac{(x-1)^2+\lambda R(x)}2.
\]
For every integer $k\equiv0\pmod {85}$, the curve
\begin{equation}\label{family}
 C_k:\quad y^2=A(x)^2-\lambda x^3(x-1)^2
\end{equation}
is smooth of genus two, its Jacobian is geometrically simple, and
\[
 \Jac(C_k)(\Q)_{\mathrm{tors}}
 \simeq\Z/2\Z\times\Z/18\Z.
\]
These curves give infinitely many geometric isomorphism classes.
\end{theorem}

\section{The rational parameter curve}
The two desired principal-divisor relations have rows
\[
 (3,2;0,0),\qquad(0,3;1,1),
\]
where the final two entries record rational Weierstrass classes.
They are realized by quadratic polynomials $A,B$ satisfying
\begin{align}
 f&=A^2-\lambda x^3(x-1)^2,\label{norm1}\\
 f&=B^2-\lambda(x-1)^3(x-r)(x-s).\label{norm2}
\end{align}
Take $A-B=(x-1)^2$ and $A+B=\lambda R(x)$. The condition
that $r$ and $s$ be Weierstrass points becomes
\[
 \lambda r^3=(r-1)^2,\qquad
 \lambda s^3=(s-1)^2.
\]
After setting $\lambda=(s-1)^2/s^3$ and removing $r=s$, the equation is
\[
 (s-1)^2r^2+s(1-2s)r+s^2=0.
\]
Its discriminant is $s^2(4s-3)$. Thus $k^2=4s-3$ gives precisely
the rational parametrization used in Theorem~\ref{thm}, with one
choice of the sign of $k$. This is an explicit genus-zero locus,
not merely a genus calculation after reduction.

\section{The torsion subgroup}
Put
\[
 B(x)=\frac{\lambda R(x)-(x-1)^2}{2},\qquad
 f(x)=A(x)^2-\lambda x^3(x-1)^2.
\]
The identities \eqref{norm1}--\eqref{norm2} hold identically, and
$B(r)=B(s)=0$. Let $W$ be the point at infinity and write
\[
 P_0=(0,A(0)),\quad P_1=(1,A(1)),\quad
 W_r=(r,0),\quad W_s=(s,0).
\]
The polynomial $R$ has degree two. Consequently $y-A$ and $y-B$
each have a pole of order five at $W$, and the norm identities give
\begin{align*}
 \divi(y-A)&=3P_0+2P_1-5W,\\
 \divi(y-B)&=3P_1+W_r+W_s-5W.
\end{align*}
In the Jacobian put $D_i=[P_i-W]$ and $T_j=[W_j-W]$. Then
\begin{equation}\label{relations}
 3D_0+2D_1=0,\qquad 3D_1+T_r+T_s=0.
\end{equation}
The distinct Weierstrass points $W,W_r,W_s$ give independent
classes $T_r,T_s$ of order two. Moreover,
\[
 A(0)=\frac{k^2+k+2}{k^2+3},\qquad A(1)=\frac\lambda2,
\]
so $P_1$ is not a Weierstrass point. Hence $2D_1\ne0$: otherwise
$2P_1\sim2W$ would force $P_1$ to be a branch point of the
unique degree-two map to $\mathbb P^1$.
The second relation gives $6D_1=0$, and the first gives
$9D_0=0$ and $3D_0=-2D_1\ne0$. Thus $D_0$ has order nine.
It follows that
\[
 \langle D_0,T_r,T_s\rangle
 \simeq\Z/9\Z\times(\Z/2\Z)^2
 \simeq\Z/2\Z\times\Z/18\Z.
\]
All these statements apply to each smooth member under consideration.

\section{Exactness and absolute simplicity}
For $k\equiv0\pmod {85}$ the denominators of the coefficients are
units at 5 and 17. The reductions of $f$ are constant on this
residue class:
\begin{align*}
 \overline f_5&=3x^5+2x^3+x^2+4x+1,\\
 \overline f_{17}&=3x^5+10x^4+16x^3+7x^2+5x+8.
\end{align*}
Both are squarefree of degree five, proving smoothness of every
member in the stated residue class. Direct counts over the prime
fields and their quadratic extensions give
\[
\begin{array}{c|rr|l|r}
 p&\#C(\F_p)&\#C(\F_{p^2})&\chi_p(X)&\#J(\F_p)\\\hline
 5&7&33&X^4+X^3+4X^2+5X+25&36\\
 17&19&321&X^4+X^3+16X^2+17X+289&324
\end{array}
\]
Prime-to-$p$ torsion injects under good reduction. The two orders
therefore bound the rational torsion order by 36, including its
5-primary and 17-primary parts. The subgroup already exhibited
has order 36, proving exactness.

For completeness, absolute simplicity can be checked without a
bound on the extension degree. The polynomial
\[
 \chi(X)=X^4+X^3+16X^2+17X+289
\]
is irreducible over $\Q$. Its Frobenius-root-ratio polynomial has
the factorization
\begin{align*}
 &\operatorname{Res}_X\bigl(\chi(X),z^4\chi(X/z)\bigr)\\
 &\qquad=17^4(z-1)^4
 (289z^4+527z^3+800z^2+527z+289)\\
 &\hspace{42mm}\cdot(17z^4+18z^3+3z^2+18z+17)^2.
\end{align*}
Both displayed quartics are irreducible over $\Q$, primitive over
$\Z$, and nonmonic; neither is cyclotomic. Therefore no ratio of
distinct Frobenius roots is a root of unity. If $\pi$ is a root of
$\chi$, then $\pi^n$ has degree four over $\Q$ for every positive
integer $n$. The reduction is thus simple over every finite
extension of $\F_{17}$, hence absolutely simple. A geometric
isogeny decomposition in characteristic zero would specialize to
one at this good prime. It follows that every $\Jac(C_k)$ in the
residue class is geometrically simple.

\section{Infinitely many moduli and a known member}
There is also a direct argument that the family is not isotrivial.
On a fixed genus-two curve there are only finitely many choices
of a Weierstrass point $W$, and for each such $W$ only finitely
many points $P$ with $18[P-W]=0$. Thus there are only finitely many
possible marked configurations $(W,P_0,P_1,W_r,W_s)$ of the kind
above on a fixed geometric curve. The coordinate on its
hyperelliptic quotient is determined by sending $W,P_0,P_1$ to
$\infty,0,1$. Consequently the coordinate $r$ can have only
finitely many values on a fixed geometric isomorphism class.
But $r=(k^2+3)/(k-1)^2$ is nonconstant and has finite fibers.
It follows that infinitely many parameters in $85\Z$ give
infinitely many geometric isomorphism classes.

At $k=0$, a small model is
\[
 y^2+(x^2+x)y=x^6+x^5-5x^4-4x^3+x^2-2x+7.
\]
It is $\Q$-isomorphic to the existing census curve $\texttt{2.18.a}$,
whose equation is
\[
 y'^2+(x'^2+x')y'=x'^6-x'^5-5x'^4+3x'^3+x'^2+2x'+7,
\]
by $x'=-x$, $y'=-x^2-y$. The new assertion is the infinitude of
geometrically simple examples with exact group $[2,18]$.

\section{Reproducibility}
The accompanying standalone Magma file
\texttt{family\_2\_18\_simple\_certificate.m} checks the generic
relations in $\Jac(C)(\Q(k))$, exact orders and independence of
the marked classes, the two reductions, the Frobenius calculations,
the root-ratio factorization, and nonconstant $G_2$-invariants.
The separate Sage file
\texttt{family\_2\_18\_independent\_check.py} verifies the polynomial
identities, counts points directly over $\F_5,\F_{25},\F_{17},\F_{289}$,
and independently checks the irreducibility and root-ratio test.
Both programs complete successfully; their logs are supplied.
\end{document}
